\subsection{Hölder 不等式}

在本节中，p 与 q 将表示满足 \(1 \leqslant p \leqslant +\infty\)、\(1 \leqslant q \leqslant +\infty\) 且以关系 \(1/p + 1/q = 1\) 相联系的两个实数；于是 \(q = p/(p - 1)\)（若 \(1 < p < +\infty\)），\(q = +\infty\)（若 p = 1），q = 1（若 \(p = +\infty\)）；p 与 q 称为共轭指数。注意，关系 \(1 \leqslant p \leqslant 2\) 等价于 \(2 \leqslant q \leqslant +\infty\)；仅当 p 与 q 都等于 2 时才有 p = q。

定理 2（Hölder 不等式）. --- 设 \(f\) 与 \(g\) 是两个几乎处处有限的数值函数，且 \(f\) 几乎处处等于 \(\mathcal{L}^p\) 中的一个函数，\(g\) 几乎处处等于 \(\mathcal{L}^q\) 中的一个函数。则函数 \(fg\)（几乎处处有定义）可积，且

\[
\mathrm{N} _ {1} (f g) \leqslant \mathrm{N} _ {p} (f) \mathrm{N} _ {q} (g).\tag{4}
\]

设 \(f_{1}\)（相应地，\(g_{1}\)）是 \(\mathcal{L}^{p}\)（相应地，\(\mathcal{L}^{q}\)）中的一个函数，\(f\)（相应地，\(g\)）几乎处处等于它；\(fg\) 几乎处处等于函数 \(f_{1}g_{1}\)，后者处处有定义且有限，且作为两个可测函数的乘积是可测的（§5, No.~3, 定理 1 的推论 5）。若 \(1 < p < +\infty\)，则关于上积分的 Hölder 不等式（Ch. I, No.~3, 命题 4）给出不等式 (4)，而关系式 \(N_{1}(fg) < +\infty\) 这时表明 \(fg\) 可积（§5, No.~6, 定理 5）。若 \(p = 1\)、\(q = +\infty\)，则不等式 (4) 以及 \(fg\) 可积这一事实是平均不等式（No.~2, 命题 1）的直接推论；于是定理在所有情形都得到证明。

推论 1. --- 设 F、G、H 是三个 Banach 空间，\((\mathbf{u}, \mathbf{v}) \mapsto \Phi(\mathbf{u}, \mathbf{v})\) 是 \(F \times G\) 到 H 中的一个连续双线性映射，使得 \(|\Phi(\mathbf{u}, \mathbf{v})| \leqslant |\mathbf{u}| \cdot |\mathbf{v}|\) 。若 \(f \in L_{F}^{p}\) 且 \(g \in L_{G}^{q}\)，则函数 \(\Phi(\mathbf{f}, \mathbf{g})\) 可积，且

\[
\left| \int \Phi (\mathbf {f}, \mathbf {g}) d \mu \right| \leqslant \int | \Phi (\mathbf {f}, \mathbf {g}) | d | \mu | \leqslant \mathrm{N} _ {p} (\mathbf {f}) \mathrm{N} _ {q} (\mathbf {g}).\tag{5}
\]

因为，\(\Phi(\mathbf{f},\mathbf{g})\) 可测（§5, No.~3, 定理 1 的推论 5）；由于 \(|\Phi(\mathbf{f},\mathbf{g})| \leqslant |\mathbf{f}| \cdot |\mathbf{g}|\)，该推论由定理 2 与 §5, No.~6 定理 5 的可积性判别法得出。

推论 1 的两个特殊情形在应用中很重要：

推论 2. --- 设 F 是一个实（相应地，复）Banach 空间，F' 是它的强对偶（TVS, III, §3, No.~1），\((\mathbf{z}, \mathbf{z}') \mapsto \langle \mathbf{z}, \mathbf{z}' \rangle\) 是 \(F \times F'\) 上的规范双线性型。若 \(f \in L_{F}^{p}\) 且 \(g \in L_{F'}^{q}\)，则实（相应地，复）函数 \(\langle f, g \rangle\) 可积，且

\[
\left| \int \langle \mathbf {f}, \mathbf {g} \rangle d \mu \right| \leqslant \int | \langle \mathbf {f}, \mathbf {g} \rangle | d | \mu | \leqslant N _ {p} (\mathbf {f}) N _ {q} (\mathbf {g}).\tag{6}
\]

因为，\(|\langle \mathbf{z},\mathbf{z}^{\prime}\rangle |\leqslant |\mathbf{z}|\cdot |\mathbf{z}^{\prime}|\)

当 F 是一个实或复 Hilbert 空间时，我们知道它可以与它的对偶 \(F'\) 规范地等同（TVS, V, §1, No.~7）。由于空间 \(L_{F}^{2}\) 完备，我们有下列结果：

推论 3. --- 设 \(\mu\) 是 X 上的一个正测度，F 是一个实（相应地，复）Hilbert 空间。在空间 \(L_{F}^{2}\) 上，对称（相应地，Hermite）型

\[
(\widetilde {\mathbf {f}}, \widetilde {\mathbf {g}}) \mapsto \int \langle \mathbf {f}, \mathbf {g} \rangle d \mu
\]

定义了一个 Hilbert 空间结构，其范数等于 \(\| \widetilde{\mathbf{f}}\| _2\)

推论 4. --- 设 F 是一个 Banach 空间，f 是 \(\mathcal{L}_{\mathrm{F}}^{p}\) 中的一个函数，g 是属于 \(\mathcal{L}^q\) 的一个数值函数；则函数 fg 可积，且

\[
\left| \int \mathbf {f} g d \mu \right| \leqslant \int | \mathbf {f} g | d | \mu | \leqslant \mathrm{N} _ {p} (\mathbf {f}) \mathrm{N} _ {q} (g).\tag{7}
\]

推论 5. --- 设 \(f_1, f_2, \ldots, f_n\) 是 \(n\) 个正的可积函数，\(\alpha_1, \alpha_2, \ldots, \alpha_n\) 是 \(n\) 个数 \(> 0\)，使得 \(\sum_{i=1}^{n} \alpha_i = 1\)；在这些条件下，函数 \(f_1^{\alpha_1} f_2^{\alpha_2} \cdots f_n^{\alpha_n}\) 可积，且

\[
\begin{array}{l} \int f _ {1} ^ {\alpha_ {1}} f _ {2} ^ {\alpha_ {2}} \dots f _ {n} ^ {\alpha_ {n}} d | \mu | \leqslant \\ \left(\int f _ {1} d | \mu |\right) ^ {\alpha_ {1}} \left(\int f _ {2} d | \mu |\right) ^ {\alpha_ {2}} \dots \left(\int f _ {n} d | \mu |\right) ^ {\alpha_ {n}}. \end{array}\tag{8}
\]

因为，乘积 \(f_{1}^{\alpha_{1}}f_{2}^{\alpha_{2}}\cdots f_{n}^{\alpha_{n}}\) 作为可测函数的乘积是可测的（§5, No.~3, 定理 1 及其推论 5）；由于不等式 (8) 对上积分成立（Ch. I, No.~2, 命题 2 的推论），函数 \(f_{1}^{\alpha_{1}}f_{2}^{\alpha_{2}}\cdots f_{n}^{\alpha_{n}}\) 可积（§5, No.~6, 定理 5），推论得证。

定理 2 的推论 2 可由下列命题加以强化：

命题 3. --- 设 \(\mu\) 是 X 上的一个正测度，F 是一个实或复 Banach 空间，F' 是它的强对偶，\((\mathbf{z},\mathbf{z}')\mapsto\langle\mathbf{z},\mathbf{z}'\rangle\) 是 F × F' 上的规范双线性型。

\(1^{\circ}\) 对每个函数 \(\mathbf{f} \in \mathcal{L}_{\mathrm{F}}^{p}\)（ \(1 \leqslant p \leqslant +\infty\)），

\[
\mathrm{N} _ {p} (\mathbf {f}) = \sup \left| \int \langle \mathbf {f}, \mathbf {g} \rangle d \mu \right|\tag{9}
\]

其中 \(\mathbf{g}\) 取遍 \(\mathcal{L}_{\mathrm{F}'}^q\) 中满足 \(\mathrm{N}_q(\mathbf{g}) \leqslant 1\) 的函数所成之集。

\(2^{\circ}\) 对每个函数 \(\mathbf{g} \in \mathcal{L}_{\mathrm{F}}^{q}\)（ \(1 \leqslant q \leqslant +\infty\)），

\[
\mathrm{N} _ {q} (\mathbf {g}) = \sup \left| \int \langle \mathbf {f}, \mathbf {g} \rangle d \mu \right|\tag{10}
\]

其中 \(\mathbf{f}\) 取遍 \(\mathcal{L}_{\mathrm{F}}^{p}\) 中满足 \(\mathrm{N}_p(\mathbf{f})\leqslant 1\) 的函数所成之集。

我们先证明关系式 (9)；分两种情形。

\begin{enumerate}
\def\labelenumi{(\roman{enumi})}
\tightlist
\item
  \(1 \leqslant p < +\infty\) 。由于当 \(\mathrm{N}_p(\mathbf{f}) = 0\) 时关系式 (9) 是平凡的（因为这时 \(\mathbf{f}\) 与 \(\langle \mathbf{f}, \mathbf{g} \rangle\) 都可忽略），我们总可以通过用标量乘 \(\mathbf{f}\) 而假设 \(\mathrm{N}_p(\mathbf{f}) = 1\) 。先设 \(\mathbf{f}\) 是一个可积阶梯函数，\(\mathbf{f} = \sum_{k=1}^{n} \dot{\mathbf{a}}_k \varphi_{\mathrm{A}_k}\)，其中诸 \(\mathrm{A}_k\) 两两不相交（§4, No.~9, 引理）。于是按假设 \(\sum_{k=1}^{n} |\mathbf{a}_k|^p \mu(\mathrm{A}_k) = 1\)。对每个 \(\varepsilon > 0\)，（对每个指标 \(k\)）存在向量 \(\mathbf{a}_k' \in \mathrm{F}'\)，使 \(|\mathbf{a}_k'|^q = |\mathbf{a}_k|^p\)（若 \(p > 1\)；相应地，\(|\mathbf{a}_k'| = 1\) 若 \(p = 1\)），且 \(\langle \mathbf{a}_k, \mathbf{a}_k' \rangle \geqslant (1 - \varepsilon) |\mathbf{a}_k| \cdot |\mathbf{a}_k'|\)（TVS, IV, §1, No.~3, 命题 8）。令 \(\mathbf{g} = \sum_{k=1}^{n} \mathbf{a}_{k}^{\prime} \varphi_{\mathrm{A}_{k}}\)，则 \(\sum_{k=1}^{n} |\mathbf{a}_{k}^{\prime}|^{q} \mu(\mathrm{A}_{k}) = 1\)（若 \(p > 1\)；相应地，\(\sup_{1 \leqslant k \leqslant n} |\mathbf{a}_{k}^{\prime}| = 1\) 若 \(p = 1\)），于是 \(\mathrm{N}_{q}(\mathbf{g}) = 1\)；另一方面，
\end{enumerate}

\[
\int \langle \mathbf {f}, \mathbf {g} \rangle d \mu = \sum_ {k = 1} ^ {n} \left\langle \mathbf {a} _ {k}, \mathbf {a} _ {k} ^ {\prime} \right\rangle \mu \left(\mathrm{A} _ {k}\right) \geqslant (1 - \varepsilon) \sum_ {k = 1} ^ {n} | \mathbf {a} _ {k} | \cdot | \mathbf {a} _ {k} ^ {\prime} | \mu (\mathrm{A} _ {k})
\]

又由于 \(|\mathbf{a}_k'| = |\mathbf{a}_k|^{p / q} = |\mathbf{a}_k|^{p - 1}\)（若 \(p > 1\)；相应地，\(|\mathbf{a}_k'| = 1\) 若 \(p = 1\)），我们有

\[
\int \langle \mathbf {f}, \mathbf {g} \rangle d \mu \geqslant (1 - \varepsilon) \sum_ {k = 1} ^ {n} | \mathbf {a} _ {k} | ^ {p} \mu (\mathrm{A} _ {k}) = (1 - \varepsilon) \mathrm{N} _ {p} (\mathbf {f}) = 1 - \varepsilon ,
\]

这就证明了此情形下的关系式 (9)。

再设 \(\mathbf{f}\) 是 \(\mathcal{L}_{\mathrm{F}}^{p}\) 中使 \(\mathrm{N}_p(\mathbf{f}) = 1\) 的任意元素。对每个 \(\varepsilon > 0\)，存在阶梯函数 \(\mathbf{f}_1 \in \mathcal{L}_{\mathrm{F}}^p\)，使 \(\mathrm{N}_p(\mathbf{f} - \mathbf{f}_1) \leqslant \varepsilon\)（§4, No.~10, 命题 19 的推论 1）。由刚才所证，存在函数 \(\mathbf{g} \in \mathcal{L}_{\mathrm{F}}^q\)，使 \(\mathrm{N}_q(\mathbf{g}) = 1\)，且

\[
\int \langle \mathbf {f} _ {1}, \mathbf {g} \rangle d \mu \geqslant \mathrm{N} _ {p} (\mathbf {f} _ {1}) - \varepsilon \geqslant 1 - 2 \varepsilon .
\]

而

\[
\int \left\langle \mathbf {f}, \mathbf {g} \right\rangle d \mu = \int \left\langle \mathbf {f} _ {1}, \mathbf {g} \right\rangle d \mu + \int \left\langle \mathbf {f} - \mathbf {f} _ {1}, \mathbf {g} \right\rangle d \mu
\]

且由 (6)，

\[
\left| \int \langle \mathbf {f} - \mathbf {f} _ {1}, \mathbf {g} \rangle d \mu \right| \leqslant \mathrm{N} _ {p} (\mathbf {f} - \mathbf {f} _ {1}) \mathrm{N} _ {q} (\mathbf {g}),
\]

由此得

\[
\left| \int \langle \mathbf {f}, \mathbf {g} \rangle d \mu \right| \geqslant 1 - 3 \varepsilon ,
\]

这就证明了 (9)。

\begin{enumerate}
\def\labelenumi{(\roman{enumi})}
\setcounter{enumi}{1}
\tightlist
\item
  \(p = +\infty\) 。我们仍可只限于 \(\mathrm{N}_{\infty}(\mathbf{f}) > 0\) 的情形。设 \(\alpha\) 是使 \(0 < \alpha < \mathrm{N}_{\infty}(\mathbf{f})\) 的任一数；按假设，由 \(x \in X\)（满足 \(|\mathbf{f}(x)| > \alpha\) 者）所成之集可测且非局部可忽略，因此它包含一个测度 \textgreater{} 0 的紧集 K。由于 f 可测，存在一个紧集 \(K_{1} \subset K\)（测度 \textgreater{} 0 的），使 f 到 \(K_{1}\) 上的限制连续。由此推出，对每个 \(\varepsilon > 0\)，存在 \(K_{1}\) 分成有限个可积集的一个分割，使 f 在每个集上的振幅 \(\leqslant \varepsilon\)；这些集中至少有一个集 A 的测度 \(>0\)。设 \(\mathbf{a}\) 是 \(\mathbf{f}\) 在 \(\mathbf{A}\) 上的一个值；则 \(|\mathbf{a}| > \alpha\)，且 \(|\mathbf{f}(x) - \mathbf{a}| \leqslant \varepsilon\) 对所有 \(x \in \mathbf{A}\) 成立。存在向量 \(\mathbf{a}' \in \mathbf{F}'\)，使 \(|\mathbf{a}'| = 1\) 且 \(|\langle \mathbf{a}, \mathbf{a}' \rangle| \geqslant |\mathbf{a}| - \varepsilon\)；函数 \(\mathbf{g} = \varphi_{\mathrm{A}} \cdot \mathbf{a}' / \mu(\mathbf{A})\) 可积且 \(\mathrm{N}_1(\mathbf{g}) = 1\)；另一方面，
\end{enumerate}

\[
\int \langle \mathbf {f}, \mathbf {g} \rangle d \mu = \frac {1}{\mu (\mathrm{A})} \int \langle \mathbf {f}, \mathbf {a} ^ {\prime} \rangle \varphi_ {\mathrm{A}} d \mu .
\]

而我们可以写出

\[
\int \langle {\bf f}, {\bf a} ^ {\prime} \rangle \varphi_ {\mathrm{A}} d \mu = \langle {\bf a}, {\bf a} ^ {\prime} \rangle \mu (\mathrm{A}) + \int \langle {\bf f} - {\bf a}, {\bf a} ^ {\prime} \rangle \varphi_ {\mathrm{A}} d \mu ,
\]

且由于

\[
\left| \langle \mathbf {f} - \mathbf {a}, \mathbf {a} ^ {\prime} \rangle \varphi_ {\mathrm{A}} \right| \leqslant \varepsilon \varphi_ {\mathrm{A}},
\]

我们看出

\[
\left| \int \langle \mathbf {f}, \mathbf {g} \rangle d \mu \right| \geqslant | \langle \mathbf {a}, \mathbf {a} ^ {\prime} \rangle | - \varepsilon \geqslant | \mathbf {a} | - 2 \varepsilon > \alpha - 2 \varepsilon ;
\]

由于 \(\varepsilon\) 任意而 \(\alpha\) 是任何 \(<\mathrm{N}_{\infty}(\mathbf{f})\) 的数，关系式 (9) 在此情形也得到验证。

证明关系式 (10) 的方式完全相同：分别考虑 \(1 \leqslant q < +\infty\) 与 \(q = +\infty\) 两种情形，并利用如下事实：对每个 \(\mathbf{z}' \in \mathrm{F}'\)，有 \(|\mathbf{z}'| = \sup_{|\mathbf{z}| \leqslant 1} |\langle \mathbf{z}, \mathbf{z}' \rangle|\)（由 \(\mathrm{F}'\) 中范数的定义）。

注. --- 1) 设 \(\mathcal{E}\) 是 \(\mathcal{L}_{\mathrm{F}'}^q\) 的一个稠密线性子空间；则当 \(\mathbf{g}\) 取遍 \(\mathcal{E}\) 与集 B——\(\mathcal{L}_{\mathrm{F}'}^q\) 中满足 \(\mathrm{N}_q(\mathbf{g}) \leqslant 1\) 的函数所成者——的交时，公式 (9) 成立。因为，只需注意到 B 的内部 \(\overset{\circ}{\mathrm{B}}\) 在 B 中稠密，且 \(\overset{\circ}{\mathrm{B}} \cap \mathcal{E}\) 在 \(\overset{\circ}{\mathrm{B}}\) 中稠密（因为 \(\overset{\circ}{\mathrm{B}}\) 是开的）。这一注记特别适用于集 \(\mathcal{E} = \mathcal{K}(\mathrm{X};\mathrm{F}')\)——由具紧支集的连续函数（取值于 \(\mathrm{F}'\)）所成者——（当 \(1 \leqslant q < +\infty\)，即 \(1 < p \leqslant +\infty\) 时）。但在这一情形，公式 (9) 当 \(\mathbf{g}\) 取遍 \(\mathrm{B} \cap \mathcal{K}(\mathrm{X};\mathrm{F}')\) 时成立，甚至对 \(p = 1\) 也如此。因为，我们可以如同上面那样只限于 \(\mathbf{f}\) 是阶梯函数的情形。这时我们看到，若 \(\mathrm{N}_1(\mathbf{f}) = 1\)，则对每个 \(\varepsilon > 0\)，存在阶梯函数 \(\mathbf{g} \in \mathcal{L}_{\mathrm{F}'}^{\infty}\)，使得 \(|\mathbf{g}(x)| \leqslant 1\) 对所有 \(x \in \mathrm{X}\) 成立，且 \(|\int \langle \mathbf{f},\mathbf{g}\rangle d\mu| \geqslant 1 - \varepsilon\) 。存在有限个两两不相交的紧集 \(\mathrm{K}_i\)，使 \(\mathbf{g}\) 取常值 \(\mathbf{a}_i'\)（在每个 \(\mathrm{K}_i\) 上），且若 K 是诸 \(\mathrm{K}_i\) 之并，则 \(\int |\mathbf{f}| \varphi_{\mathbf{C} K} d\mu \leqslant \varepsilon\) 。设 \(\mathrm{U}_i\) 是 \(\mathrm{K}_i\) 的一个邻域，使诸 \(\mathrm{U}_i\) 两两不相交，并设 \(h_i\) 是 X 到 {[}0,1{]} 中的一个连续映射，其支包含于 \(\mathrm{U}_i\) 且在 \(\mathrm{K}_i\) 上等于 1。令 \(\mathbf{h} = \sum \mathbf{a}_i' h_i\)，则在 K 上 \(\mathbf{h}(x) = \mathbf{g}(x)\)，且在 X 上 \(|\mathbf{h}(x)| \leqslant 1\)，因此

\[
\int | \langle \mathbf {f}, \mathbf {h} \rangle | \varphi_ {\mathbf {C K}} d \mu \leqslant \varepsilon
\]

从而 \(|\int \langle \mathbf{f},\mathbf{h}\rangle d\mu |\geqslant 1 - 3\varepsilon\)，这就证明了我们的断言。对公式 (10) 可作类似的注记。

\begin{enumerate}
\def\labelenumi{\arabic{enumi})}
\setcounter{enumi}{1}
\tightlist
\item
  设 \(\mu\) 是 \(X\) 上的一个正测度，\(f\) 是一个可测函数 \(\geqslant 0\)（有限或无限），其支集含于紧集 \(K_{n}\) 的一个可数并。则对每个 \(p\)（满足 \(1 \leqslant p \leqslant +\infty\) 者），
\end{enumerate}

\[
\mathrm{N} _ {p} (f) = \sup \int^ {*} | f g | d \mu ,\tag{11}
\]

其中 \(g\) 取遍 \(\mathcal{K}(\mathrm{X};\mathbf{R})\) 中满足 \(\mathrm{N}_q(g)\leqslant 1\) 的函数所成之集。因为，当 \(\mathrm{N}_p(f) < +\infty\) 时公式 (11) 是 (9) 的特殊情形，因为这时 \(f\) 等价于 \(\mathcal{L}^p\) 中的一个函数（§5, No.~6, 定理 5）。若 \(\mathrm{N}_p(f) = +\infty\)，对每个整数 \(n > 0\) 令 \(f_{n} = \inf (n,f\varphi_{\mathrm{K}_{n}})\) 。则

\[
\mathrm{N} _ {p} (f _ {n}) = \sup \int^ {*} | f _ {n} g | d \mu \leqslant \sup \int^ {*} | f g | d \mu ,
\]

于是（假设如我们所可以的那样，序列 \((\mathbf{K}_n)\) 递增）过渡到极限，便得 \(\sup \int^{*}|fg| d\mu = +\infty\)（§1, No.~3, 定理 3）。

推论. --- 设 \(\mu\) 是 X 上的一个正测度，F 是一个 Banach 空间，\(F'\) 是它的强对偶，g 是 \(L_{F'}^{q}\) 中的任一函数。则由 \(L_{F}^{p}\) 上的线性形式 \(f \mapsto \int \langle f, g \rangle d\mu\) 过渡到商得到的 \(L_{F}^{p}\) 上的线性形式是连续的，且范数为 \(N_{q}(g)\) 。

